% chapters/02-sources.tex -- 세 글: 제목, URL, 핵심, 서로의 관계.
% ONLY source: perf.sources.md (S1-S3, fetched 2026-10-14) and perf.audit.md §0.
% Summaries were extracted by WebFetch and are not checked against the Next.js docs.
\chapter{세 글}\label{ch:sources}

% src: perf.sources.md header, perf.audit.md §0
\paragraph{결과.}
세 글은 한 사슬이다.
1편은 이동이 왜 멈추는지, 2편은 읽기를 어디서 하는지, 3편은 읽기를 어떤 순서로 하는지를 말한다.
\Cref{fig:src-chain}이 세 글의 관계를 보인다.
글 요약은 WebFetch로 뽑은 것이다. Next.js 공식 문서와 대조하지 않았다.

\begin{inbrief}
  \item[Why] 점검 25개의 기준이 이 세 글이다. 기준을 먼저 알아야 항목이 읽힌다.
  \item[When] 03--06장의 항목을 읽기 전이다.
  \item[Where] 저자 \texttt{k-svelte-master}가 velog에 쓴 ``Next.js 이렇게 쓰면 느립니다'' 연재다.
  \item[How] 글마다 제목, URL, 핵심 5개를 적는다. 마지막에 관계 그림을 둔다.
\end{inbrief}

\section{S1. 1편: 동적 RSC와 prefetch}\label{sec:src-s1}

\paragraph{무엇을 알려고 쓰나.}
이동이 느린 첫 번째 이유는 무엇이고, 무엇으로 고치는가.

\begin{description}[font=\bfseries\small, leftmargin=3.4em, labelwidth=3em, itemsep=1pt]
  \item[제목] Next.js 이렇게 쓰면 느립니다 1편: 동적 RSC(dynamic RSC)와 prefetch
  \item[URL] {\footnotesize\url{https://velog.io/@k-svelte-master/nextjs-slow-1-dynamic-rsc-prefetch}}
\end{description}

% src: perf.sources.md S1
\begin{enumerate}[nosep]
  \item 이동은 RSC 구멍(hole)의 payload가 올 때까지 멈춘다. React는 덜 된 tree를 그릴 수 없다.
  \item 요청 시점 API(request-time API)가 route 어디에든 있으면 route 전체가 동적(dynamic)이 된다. 예는 \texttt{cookies()}, \texttt{headers()}, \texttt{searchParams}, 이를 쓰는 auth helper다.
  \item 동적 route는 page payload를 prefetch하지 않는다. \texttt{loading.js}까지만 받는다. 클라이언트 라우터 캐시(client router cache)는 0초 보존이라 클릭마다 왕복(round trip)이 생긴다.
  \item 처방은 \texttt{cacheComponents: true}다. 동적 읽기를 \texttt{<Suspense>}로 감싸 static shell이 prefetch와 cache에 오르게 한다. \texttt{params}는 잎에서만 읽고, 자주 여는 segment는 \texttt{generateStaticParams}를 쓰고, \texttt{experimental.staleTimes}를 일부러 정한다.
  \item 규칙 하나: 이동 속도는 클릭 전에 RSC 구멍이 prefetch와 cache로 채워졌는지로 정해진다.
\end{enumerate}

\section{S2. 2편: Server Actions는 순차 실행된다}\label{sec:src-s2}

\paragraph{무엇을 알려고 쓰나.}
읽기를 Server Action으로 하면 왜 안 되는가.

\begin{description}[font=\bfseries\small, leftmargin=3.4em, labelwidth=3em, itemsep=1pt]
  \item[제목] Next.js 이렇게 쓰면 느립니다 2편: Server Functions(Server Actions)는 순차 실행된다
  \item[URL] {\footnotesize\url{https://velog.io/@k-svelte-master/nextjs-slow-2-server-actions-sequential}}
\end{description}

% src: perf.sources.md S2
\begin{enumerate}[nosep]
  \item Server Action은 client마다 하나씩 줄을 선다. 1초짜리 action 셋을 \texttt{Promise.all}로 불러도 약 3초다.
  \item 첫 렌더(initial render) 중에는 부를 수 없다. TanStack Query는 \texttt{queryFn}으로 쓰기를 권하지 않고, \texttt{useSuspenseQuery}는 실패한다.
  \item 설계 규칙: Server Action은 변경(mutation) 전용이다. 읽기는 Server Component나 Route Handler로 한다.
  \item BFF(backend for frontend): backend 호출을 Next server를 거쳐 한다. 같은 곳에 있어 왕복이 약 3 ms다. 브라우저에서 backend로 가면 왕복마다 약 100 ms다.
\end{enumerate}

\section{S3. 3편: Suspense, 스트리밍, 워터폴}\label{sec:src-s3}

\paragraph{무엇을 알려고 쓰나.}
같은 읽기를 어떤 순서로 하면 시간이 줄어드는가.

\begin{description}[font=\bfseries\small, leftmargin=3.4em, labelwidth=3em, itemsep=1pt]
  \item[제목] Next.js 이렇게 쓰면 느립니다 3편: Suspense, 스트리밍(streaming), 워터폴(waterfall)
  \item[URL] {\footnotesize\url{https://velog.io/@k-svelte-master/nextjs-slow-3-suspense-streaming-waterfall}}
\end{description}

% src: perf.sources.md S3
\begin{enumerate}[nosep]
  \item Suspense 스트리밍: 첫 chunk에 fallback이 가고, 내용은 뒤 chunk로 온다. 마지막 경계가 끝나야 응답이 닫힌다.
  \item 중첩 async 컴포넌트는 워터폴을 만든다. 자식의 fetch는 부모의 await 뒤에 시작해서 300 ms + 300 ms = 600 ms다.
  \item 처방: \texttt{React.cache}로 감싼다. 부모에서 \texttt{void getComments(id)}로 미리 시작(preload)하고 자식은 같은 promise를 기다린다. 시간은 약 300 ms다.
  \item client fetch(\texttt{useQuery})만 쓰면 응답이 loading 상태로 끝난다. 크롤러(crawler)는 덜 된 HTML을 받는다. 색인할 내용은 Suspense와 server fetch로 낸다.
  \item TanStack hydration: server 값은 props로만 client에 닿는다. \texttt{prefetchQuery}와 \texttt{HydrationBoundary} 없이 쓰면 요청이 두 번 나고 hydration 불일치가 생긴다.
\end{enumerate}

\section{세 글의 관계}\label{sec:src-chain}

\paragraph{무엇을 알려고 쓰나.}
세 글이 한 이동에서 어떤 순서로 만나는가.

\begin{figure}[H]
  \centering
  \resizebox{\linewidth}{!}{%
  \begin{tikzpicture}[font=\scriptsize, node distance=1.0cm]
    \node[flowbox, minimum width=3.4cm] (click) {클릭};
    \node[badbox, minimum width=3.6cm, right=1.2cm of click] (s1)
      {\textbf{S1} 동적 route\\prefetch 없음, 캐시 0초\\왕복이 매번 생김};
    \node[flowbox, minimum width=3.6cm, right=1.2cm of s1] (s2)
      {\textbf{S2} 읽기는 어디서\\action은 변경 전용\\읽기는 RSC나 Route Handler};
    \node[flowbox, minimum width=3.6cm, right=1.2cm of s2] (s3)
      {\textbf{S3} 읽기를 어떤 순서로\\Suspense, preload\\600 ms를 300 ms로};
    \node[goodbox, minimum width=3.0cm, right=1.2cm of s3] (done) {화면과 색인 HTML};
    \draw[flowarrow] (click) -- (s1);
    \draw[flowarrow] (s1) -- (s2) node[lbl, midway, above] {구멍을\\채울 읽기};
    \draw[flowarrow] (s2) -- (s3) node[lbl, midway, above] {읽기가\\여럿이면};
    \draw[flowarrow] (s3) -- (done);
    \node[lbl, below=0.5cm of s1, text width=3.6cm] (n1) {원인: 요청 시점 API};
    \node[lbl, below=0.5cm of s2, text width=3.6cm] (n2) {원인: 줄 서는 action};
    \node[lbl, below=0.5cm of s3, text width=3.6cm] (n3) {원인: 중첩 await};
  \end{tikzpicture}}
  \par\smallskip{\centering\scriptsize \flowkey{badbox} 원인 또는 문제 \quad \flowkey{goodbox} 결과 \quad \flowkey{flowbox} 일반 단계\par}
  \caption{세 글의 관계. 1편이 이동을 멈추고, 2편과 3편이 그 멈춤 안의 읽기를 다룬다.}
  \label{fig:src-chain}
  % src: perf.sources.md S1-S3 (the arrows are the book's reading of the three articles)
\end{figure}

\paragraph{읽는 법.}
\begin{itemize}[nosep]
  \item S1은 이동마다 왕복이 생기는 이유를 다룬다.
  \item S2는 읽기를 둘 자리를 정한다. 그 자리가 Server Component나 Route Handler다.
  \item S3는 그 자리 안에서 읽기를 겹치는 법이다. 순서가 틀리면 시간이 더해진다.
\end{itemize}
